\begin{textsample}{POS Dim 4 – 2009\_c6 – Score -3.00 – 2009\_c6/enade\_2009\_tec\_des\_mod\_3927bb5ec2.txt}  \label{ex:f4_pos_324}
Capítulo 3 Análise Técnica da Prova Este capítulo tem por objetivo apresentar o desempenho dos estudantes de TECNOLOGIA EM DESIGN DE MODA no ENADE/2009.
Para isso, foram calculadas as estatísticas básicas da prova como um todo, bem como as estatísticas das questões relacionadas à Formação Geral e ao Componente Específico.
Nas Tabelas, são apresentadas as seguintes estatísticas : total da população e de presentes, média, erro-padrão da média, desvio-padrão, mediana, nota máxima, nota mínima e coeficiente de assimetria ( ver definição em 1.4.3 ).
Nas estatísticas das questões discursivas, foram detalhadas também as notas nulas : em branco, por protesto ou por insuficiência.
As estatísticas apresentadas neste capítulo contemplam, separadamente, os ingressantes, os concluintes e o total de estudantes, e foram calculadas tendo -se em vista as seguintes agregações : ( a ) Região Geográfica e País e ( b ) a Categoria Administrativa e a Organização Acadêmica.
As notas brutas das questões discursivas em cada uma das partes, Formação Geral/Componente Específico, foram calculadas a partir da média simples das notas de cada uma das questões correspondentes.
A primeira questão discursiva de Formação Geral foi desconsiderada e, portanto, a nota bruta referente às questões discursivas de Formação Geral foi igual à nota da questão mantida, a segunda.
No componente específico, a nota bruta correspondeu à média das notas nas três questões.
As notas brutas das questões objetivas em cada uma das partes, Formação Geral/Componente Específico, foram calculadas a partir dos acertos dos estudantes considerando -se uma regra de três para a conversão para a escala de 0 a 100 descontando -se aquelas questões que foram desconsideradas pela comissão ou que não passaram pelo teste do índice de discriminação ponto-bisserial ( ver seção 1.4.6 para a descrição do índice ).
Questões com índices de discriminação inferior a 0,2 foram excluídas do cômputo da nota do estudante.
As notas brutas de Formação Geral e do Componente Específico para cada um dos estudantes foram calculadas como médias ponderadas das suas notas nas questões objetivas e discursivas, ponderação diferente para cada parte : C\_FG = 0,6 × C\_FG^OBJ + 0,4 × C\_FG^DIS C\_CE = 0,85 × C\_CE^OBJ + 0,15 × C\_CE^DIS, onde C\_FG é a nota bruta da Formação Geral, C\_FG^OBJ é a nota bruta nas questões objetivas da Formação Geral, C\_FG^DIS é a nota bruta nas questões discursivas da Formação Geral, C\_CE é a nota bruta do Componente Específico, C\_OBJ CE é a nota bruta nas questões objetivas do Componente Específico, e C\_DIS CE é a nota bruta nas questões discursivas do Componente Específico.
A nota bruta da prova, C\_GER, por sua vez, é também uma média ponderada da nota de cada uma das partes, Formação geral/Componente Específico : C\_GER = 0,25 x C\_FG + 0,75 x C\_CE. 3.1 Estatísticas Básicas da Prova 3.1.1 Estatísticas Básicas Gerais A Tabela 3.1 apresenta as estatísticas básicas da prova por grupo de estudantes.
O total de estudantes inscritos para a prova foi de 4.820.
Destes, 78,2 \% compareceram, sendo que a presença foi maior entre concluintes ( 87,1 \% ) do que entre ingressantes ( 74,3 \% ).
A média geral da prova foi de 46,51, sendo que os ingressantes obtiveram média mais baixa, de 44,31, que os concluintes, cuja média foi 50,74.
O desvio-padrão geral foi de 14,1, sendo que o do grupo de ingressantes foi igual ao do grupo de concluintes ( 13,8 ), não indicando uma maior homogeneidade nas notas para nenhum dos grupos.
A nota máxima foi 89,3, obtida por pelo menos um concluinte ao passo que a maior nota obtida por um ingressante foi 85,0.
Os coeficientes de assimetria para a população e para o grupo de concluintes foram negativos, indicando uma maior concentração à direita da distribuição e um maior espalhamento à esquerda.
Por outro lado, para o grupo de ingressantes o coeficiente de assimetria foi nulo, indicando uma distribuição equilibrada sem concentração e/ou espalhamento.
Uma comparação entre o desempenho de ingressantes e concluintes pode ser realizada a partir dos dados apresentados no Gráfico 3.1.
Cumpre destacar que as notas dos concluintes, de um modo geral, foram mais elevadas do que as dos ingressantes, cuja distribuição de notas está à esquerda dos primeiros.
A ideia de melhor desempenho é reforçada pelas notas medianas dos dois grupos : 51,6 entre os concluintes e 44,8 entre os ingressantes.
Nas faixas de notas até 50 pontos o percentual de ingressantes é maior que o dos concluintes, situação esta que se inverte para as faixas de notas superiores a 50 pontos.
Assim, aproximadamente 35,1 \% dos ingressantes situaram -se nas faixas de notas superiores a 50,0 pontos, por oposição a 53,5 \% dos concluintes.
O Gráfico 3.2 apresenta informações referentes ao desempenho geral de ingressantes e concluintes, comparando os resultados considerando a Região, Categoria Administrativa e Organização Acadêmica.
Levando -se em conta as notas médias dos estudantes em cada Região, observa -se, nos dois grupos, uma variação expressiva entre as Regiões.
As notas médias mais elevadas entre os concluintes ocorreram nas Regiões Sul ( 52,45 ) e Centro-Oeste ( 51,73 ), ambas acima da média nacional ( 50,73 ).
Em relação aos ingressantes, as médias mais altas foram obtidas nas Regiões Norte ( 48,27 ), Centro-Oeste ( 46,23 ) e Sul ( 45,96 ), nas demais Regiões a média foi inferior à média nacional ( 44,31 ).
Na Região Nordeste foi obtida a média de notas mais baixa entre os concluintes ( 49,68 ) e na Região Sudeste entre os ingressantes ( 42,69 ).
Cursos nesta área foram oferecidos em instituições federais, estaduais, municipais e privadas.
Levando -se em conta os agrupamentos dos estudantes em categorias administrativas, observa -se que a pontuação média mais elevada entre os concluintes foi encontrada em instituições de origem federal ( 58,69 ) – acima da média nacional.
Também se observa que a pontuação média mais elevada entre os ingressantes foi encontrada em instituições de origem estadual ( 53,70 ) – também acima da média nacional.
Tendo como foco as notas médias dos estudantes agrupados por Organização Acadêmica, notamos que as faculdades apresentaram as maiores médias ( concluintes – 50,98 e ingressantes – 45,61 ), ao passo que as piores notas médias foram observadas nos Centros Universitários. 3.1.2 Estatísticas Básicas de Formação Geral A Tabela 3.2 apresenta as estatísticas básicas de Formação Geral por grupo de estudantes.
O total de estudantes inscritos no exame foi de 4.820.
Destes, 78,2 \% compareceram, sendo que a presença foi maior entre concluintes ( 87,1 \% ) do que entre ingressantes ( 74,3 \% ).
A média geral da parte da prova de Formação Geral foi de 42,9, sendo que os ingressantes obtiveram média mais baixa, 42,15, que os concluintes, cuja média foi 44,22.
O desvio-padrão geral foi de 21,0, sendo o do grupo de ingressantes menor ( 20,9 ), do que o do grupo dos concluintes ( 21,1 ), aqueles possuíram uma distribuição mais homogênea nas notas.
A nota máxima foi 100,0, obtida por estudantes de ambos os grupos.
Para esta parte, assim como para a nota geral, os coeficientes de assimetria são positivos para todos os grupos.
Como, porém, os valores são bem perto de zero, estes indicaram uma ligeira inclinação para a esquerda e uma cauda mais comprida para a direita.
O Gráfico 3.3 permite a comparação entre concluintes e ingressantes relativamente ao desempenho na parte da prova que aborda a Formação Geral.
Nesta parte, tal como na prova como um todo, também os concluintes apresentaram desempenho superior, embora seja perceptível um maior equilíbrio.
A mediana dos concluintes foi um pouco maior do que a dos ingressantes, o mesmo acontecendo com a média.
Dos ingressantes, aproximadamente 36 \% alcançaram nota superior a 50 ; dentre os concluintes, o percentual dos estudantes nesta faixa de notas foi um pouco maior, 40 \%.
No Gráfico 3.4 são apresentadas as informações referentes ao desempenho dos ingressantes e dos concluintes na parte de Formação Geral, em diferentes agregações como Região, Categoria Administrativa e Organização Acadêmica.
Considerando -se as notas médias dos estudantes segundo cada Região, observou -se que, tanto no grupo dos ingressantes quanto no grupo dos concluintes, houve diferenças marcantes.
As médias variam, entre os concluintes, de 37,61 na Região Norte a 46,72 na Região Sul.
Entre os ingressantes, os valores foram de 36,77 na Região Nordeste a 44,22 na Região Sul.
A maior dispersão de notas dos estudantes de uma dada Região foi encontrada na Norte tanto para os concluintes quanto para os ingressantes ( respectivamente 9,35 e 7,30 ).
Nas Regiões Nordeste, Sudeste e Sul os concluintes apresentaram médias maiores do que os ingressantes, a situação inversa aconteceu no Norte e no Centro-Oeste.
Todas as Regiões apresentaram concluintes.
Tendo como foco as categorias administrativas, destacou -se que todas as categorias estavam representadas, sendo as notas médias mais baixas encontradas nas instituições estaduais, tanto para concluintes ( 34,90 ) quanto para ingressantes ( 42,68 ) e as mais altas nas instituições federais ( respectivamente 55,84 e 49,30 ).
Note que, com este recorte, com exceção das instituições estaduais, concluintes apresentaram médias de notas superiores aos dos ingressantes, semelhantemente ao ocorrido para o País como um todo.
Os estudantes das instituições privadas se apresentam mais homogêneos do que os das outras categorias, já que o desvio padrão para concluintes e ingressantes é o menor.
O Considerando -se o tipo de Organização Acadêmica, os Centros Universitários apresentam as maiores médias entre os concluintes ( 46,09 ), por outro lado as Universidades apresentam as maiores médias entre os ingressantes ( 43,55 ).
As Faculdades apresentaram as menores médias para concluintes ( 42,49 ) e Centro Universitários para ingressantes ( 40,49 ).
Os Centros Federais de Educação Tecnológica e os Institutos Federais de Educação, Ciência e Tecnologia não estavam representadas nesta área. 3.1.3 Estatísticas Básicas do Componente Específico A Tabela 3.3 apresenta as estatísticas básicas referentes ao Componente Específico da prova, tendo sua média geral, de 47,78, superior ao desempenho da prova de TECNOLOGIA EM DESIGN DE MODA como um todo, cuja média foi 46,5, como apresentado na Tabela 3.1, que por sua vez foi superior ao desempenho na Formação Geral, que teve média igual a 42,9, conforme mostrado na Tabela 3.2.
Os concluintes apresentaram um desempenho médio ( 53,0 ) superior ao dos ingressantes ( 45,1 ), como ocorrido para a nota da Formação geral, havendo ainda indícios de homogeneidade entre as notas médias dos ingressantes e dos concluintes, como sugeriu o desvio padrão, praticamente igual para o grupo dos ingressantes e dos concluintes.
As notas máximas obtidas por concluintes e ingressantes foram iguais a, 88,1 e 84,0 ; respectivamente.
As distribuições de notas para cada grupo se apresentaram com uma concentração nas notas maiores e uma cauda mais comprida para as notas menores, como evidenciados pelos coeficientes de assimetria negativos.
A situação de assimetria é confirmada pelo Gráfico 3.5, que proporciona uma comparação do desempenho entre concluintes e ingressantes em relação ao Componente Específico da prova.
Claramente a distribuição das notas médias dos ingressantes, à direita daquela dos concluintes, evidenciou o percentual mais alto de notas médias nas faixas mais baixas para este grupo.
As medianas das distribuições das notas dos ingressantes e dos concluintes foram 45,4 e 54,4, respectivamente, enfatizando o melhor desempenho dos concluintes nesta parte da prova.
Assim como os Gráficos 3.2 e 3.4, o Gráfico 3.6 apresenta uma comparação dos resultados em relação às Regiões, Categoria Administrativa e Organização Acadêmica, agora levando em conta o desempenho dos ingressantes e dos concluintes no Componente Específico da prova.
Tendo como foco as notas médias dos estudantes, em determinada Região, observa -se que, entre os concluintes, as mais elevadas foram encontradas nas Regiões Norte ( 54,78 ), Centro-Oeste ( 54,56 ) e Nordeste ( 53,45 ), valores acima da média nacional para esta parte ( 52,95 ).
Entre os ingressantes, as maiores notas ocorreram também nas mesmas Regiões : Norte ( 50,18 ), Centro-Oeste ( 47,05 ) e Nordeste ( 46,62 ), todas acima da média nacional ( 45,08 ).
Novamente, para a análise segundo a Categoria Administrativa, todos os tipos de entidades foram representadas.
As médias mais elevadas entre os concluintes foram encontradas em instituições de origem federal ( 59,68 ), entre os ingressantes, foram encontradas nas estaduais ( 57,41 ).
As médias mais baixas ocorreram entre os concluintes nas instituições municipais ( 50,14 ) e entre os ingressantes nas instituições particulares ( 44,55 ).
Quanto à Organização Acadêmica, o comentário é o mesmo já feito para a nota bruta e referente à Formação Geral : somente Universidades, Centros Universitários e Faculdades foram representadas com este curso.
Considerando -se o tipo de Organização Acadêmica, a ordenação das médias é a mesma para concluintes e ingressantes : Faculdades ( 53,86 e 47,09 respectivamente ) seguidas das Universidades ( 52,74 e 45,66 respectivamente ) e dos Centros Universitários ( 52,09 e 42,13 respectivamente ). 3.2 Análise das Questões Objetivas 3.2.1 Formação Geral A Tabela 3.4 apresenta as estatísticas básicas relativas às oito questões objetivas da Formação Geral dos estudantes.
Adiantamos que a primeira questão foi desconsiderada pelo critério da correlação ponto-bisserial.
Como pode ser observado na Tabela 3.4, os estudantes concluintes obtiveram desempenho superior, com nota média de 53,9, em relação aos estudantes ingressantes cuja nota média foi 52,5.
A variabilidade entre os concluintes foi um pouco menor do que entre os ingressantes, como indicaram os desvios-padrão de 22,7, para o grupo de ingressantes e 22,1 para o grupo dos concluintes.
A amplitude total das notas foi a mesma para ambos os grupos : nota mínima de zero e máxima de 100,0.
A mediana foi igual para ambos os grupos : 57,1.
A assimetria, medida pelo coeficiente de assimetria ( skewness ) foi nula para os ingressantes ( 0,0 ) e negativa para os concluintes ( -0,1 ).
O grupo de ingressantes apresenta simetria mesmo com a média inferior à mediana.
O grupo de concluintes apresenta assimetria com concentração à direita com cauda mais pesada para a esquerda, impressão reforçada, também, por uma média de valor inferior à mediana correspondente.
A Tabela 3.5 apresenta a classificação das questões objetivas da Formação Geral, segundo o índice de facilidade ( questões em vermelho foram desconsideradas pelo critério ponto-bisserial ).
Como as questões foram comuns a todas as áreas, a tabela apresenta a classificação considerando -se os concluintes de todas as áreas e aqueles do curso de TECNOLOGIA EM DESIGN DE MODA.
De acordo com os índices obtidos, as questões objetivas da prova foram assim avaliadas para a área de TECNOLOGIA EM DESIGN DE A Tabela 3.6 descreve os resultados das análises das questões objetivas relativas à Formação Geral, segundo o poder de discriminação, utilizando -se para tal o índice de discriminação ponto-bisserial.
Como as questões foram comuns a todas as áreas, a tabela apresenta a classificação considerando -se os concluintes de todas as áreas e aqueles do curso de TECNOLOGIA EM DESIGN DE MODA.
A desconsideração da questão segundo este critério considerou o universo de todos os concluintes.
Nenhuma questão apresentou índice igual ou acima de 0,40 ( nenhuma para a população de concluintes como um todo ), sendo, portanto, classificadas como muito boas nesta característica, para este grupo de estudantes, três questões foram classificadas como boas com índice entre 0,30, inclusive, e 0,40 ( também três para os concluintes como um todo ) e outras três foram classificadas como médias, com índice entre 0,20, inclusive, e 0,30 ( quatro para os concluintes como um todo ).
A primeira questão foi classificada como fraca tanto para a população de concluintes como um todo como para a de TECNOLOGIA EM DESIGN DE MODA e foi desconsiderada pelo critério de discriminação bisserial.
A quarta questão foi classificada como fraca com respeito ao índice de discriminação somente na área e não foi desconsiderada.
MODA : das oito questões, duas tiveram entre 60 \%, inclusive, e 85 \% de acertos e foram classificadas como fáceis ( somente uma para o total de concluintes ) ; outras três foram consideradas medianas, alcançando entre 40 \%, inclusive, e 60 \% de acertos ( também três para o total de concluintes ) ; duas obtiveram entre 15 \%, inclusive, e 40 \% de acertos, sendo consideradas difíceis ( três para o total de concluintes ).
Por fim, uma foi classificada como muito difícil tendo obtido percentual menor do que 15 \% de acertos ( uma para o total de concluintes ), coincidentemente a questão desconsiderada pelo critério ponto-bisserial.
Nenhuma questão foi classificada na categoria Muito fácil.
A Tabela 3.6 descreve os resultados das análises das questões objetivas relativas à Formação Geral, segundo o poder de discriminação, utilizando -se para tal o índice de discriminação ponto-bisserial.
Como as questões foram comuns a todas as áreas, a tabela apresenta a classificação considerando -se os concluintes de todas as áreas e aqueles do curso de TECNOLOGIA EM DESIGN DE MODA.
A desconsideração da questão segundo este critério considerou o universo de todos os concluintes.
Nenhuma questão apresentou índice igual ou acima de 0,40 ( nenhuma para a população de concluintes como um todo ), sendo, portanto, classificadas como muito boas nesta característica, para este grupo de estudantes, três questões foram classificadas como boas com índice entre 0,30, inclusive, e 0,40 ( também três para os concluintes como um todo ) e outras três foram classificadas como médias, com índice entre 0,20, inclusive, e 0,30 ( quatro para os concluintes como um todo ).
A primeira questão foi classificada como fraca tanto para a população de concluintes como um todo como para a de TECNOLOGIA EM DESIGN DE MODA e foi desconsiderada pelo critério de discriminação bisserial.
A quarta questão foi classificada como fraca com respeito ao índice de discriminação somente na área e não foi desconsiderada.
Na sequência, no Quadro 3.1, estão especificadas, por questão, as habilidades aferidas pelas questões de Formação Geral da prova da TECNOLOGIA EM DESIGN DE MODA, bem como os resultados da análise relativa aos índices de facilidade e de discriminação de cada questão.
O índice de facilidade para a área de TECNOLOGIA EM DESIGN DE MODA teve uma variação considerável, oscilando entre 0,08 até 0,81 ( na população de concluintes o intervalo foi de 0,08 a 0,75 ) e o de discriminação ficou entre 0,03 e 0,39 ( na população de concluintes o intervalo foi de 0,02 a 0,39 ).
A questão 7 aferiu as habilidades de “ Analisar e criticar informações ”, “ Extrair conclusões por indução e/ou dedução ”, “ Estabelecer relações, comparações e contrastes em diferentes situações ”, “ Detectar contradições ”, “ Fazer escolhas valorativas avaliando consequências ”, “ Questionar a realidade ”, e “ Argumentar coerentemente ” foi a que apresentou, para a população de concluintes de TECNOLOGIA EM DESIGN DE MODA, maior poder discriminatório com índice de 0,39.
A questão mais fácil, com 81 \% de acertos, foi a de número 8, que aferiu também, além das já listadas, a habilidade de “ Ler e interpretar textos ”.
De outra parte, a questão com menor índice de facilidade e de discriminação em Formação Geral foi a questão 1, que foi desconsiderada, e que deveria aferir as habilidades de “ Ler e interpretar textos ”, “ Analisar e criticar informações ”, “ Extrair conclusões por indução e/ou dedução ”, “ Estabelecer relações, comparações e contrastes em diferentes situações ”, “ Detectar contradições ”, “ Fazer escolhas valorativas avaliando conseqüências ”, “ Questionar a realidade ”, e “ Argumentar coerentemente ” Na sequência, no Quadro 3.1, estão especificadas, por questão, as habilidades aferidas pelas questões de Formação Geral da prova da TECNOLOGIA EM DESIGN DE MODA, bem como os resultados da análise relativa aos índices de facilidade e de discriminação de cada questão.
O índice de facilidade para a área de TECNOLOGIA EM DESIGN DE MODA teve uma variação considerável, oscilando entre 0,08 até 0,81 ( na população de concluintes o intervalo foi de 0,08 a 0,75 ) e o de discriminação ficou entre 0,03 e 0,39 ( na população de concluintes o intervalo foi de 0,02 a 0,39 ).
A questão 7 aferiu as habilidades de “ Analisar e criticar informações ”, “ Extrair conclusões por indução e/ou dedução ”, “ Estabelecer relações, comparações e contrastes em diferentes situações ”, “ Detectar contradições ”, “ Fazer escolhas valorativas avaliando consequências ”, “ Questionar a realidade ”, e “ Argumentar coerentemente ” foi a que apresentou, para a população de concluintes de TECNOLOGIA EM DESIGN DE MODA, maior poder discriminatório com índice de 0,39.
A questão mais fácil, com 81 \% de acertos, foi a de número 8, que aferiu também, além das já listadas, a habilidade de “ Ler e interpretar textos ”.
De outra parte, a questão com menor índice de facilidade e de discriminação em Formação Geral foi a questão 1, que foi desconsiderada, e que deveria aferir as habilidades de “ Ler e interpretar textos ”, “ Analisar e criticar informações ”, “ Extrair conclusões por indução e/ou dedução ”, “ Estabelecer relações, comparações e contrastes em diferentes situações ”, “ Detectar contradições ”, “ Fazer escolhas valorativas avaliando conseqüências ”, “ Questionar a realidade ”, e “ Argumentar coerentemente ” O Gráfico 3.7 mostra, a título de ilustração, o comportamento das respostas à questão 7 de Formação Geral para os estudantes concluintes de TECNOLOGIA EM DESIGN DE MODA como função da nota geral da prova4.
O poder de discriminação foi bom tanto para os estudantes da área como também, para o conjunto de todos os concluintes.
Neste Gráfico, cada uma das cinco curvas representou o percentual de respostas em determinada alternativa da questão, em função do escore total dos estudantes concluintes da área de TECNOLOGIA EM DESIGN DE MODA.
A curva azul marinho corresponde à alternativa A, a correta para esta questão.
Assim, observa -se que entre os estudantes com notas mais baixas – desempenho geral mais fraco na prova como um todo –, a tendência foi a escolha de alternativas incorretas, especialmente a alternativa B – curva verde claro e da alternativa D – curva verde escuro.
A proporção de estudantes que selecionaram a resposta correta A, aumentou com o desempenho deste grupo de estudantes chegando ao acerto de 100 \% dos estudantes com notas pelo menos 87,5.
Esta análise permitiu verificar como a questão discriminou entre os grupos de desempenho.
Os Gráficos relativos às demais questões de Formação Geral constam do Anexo I. 3.2.2 Componente Específico A Tabela 3.7 apresenta as estatísticas básicas em relação às questões objetivas do Componente Específico.
Como pode ser observado, os estudantes concluintes obtiveram desempenho superior, mas por uma pequena margem com uma média de 59,47 pontos, em relação aos estudantes ingressantes com uma média de 51,46 pontos.
Os desvios-padrão indicaram que a variabilidade entre os concluintes foi igual a 15,9 e entre os ingressantes, 16,2.
Os coeficientes de assimetria para os dois grupos foram negativos indicando uma inclinação para a direita.
A assimetria foi, porém, maior entre os concluintes.
A análise das questões objetivas do Componente Específico em TECNOLOGIA EM DESIGN DE MODA, em relação ao grau de facilidade da prova ( ver Tabela 3.8 ), mostrou que, das 27 questões, onze encontraram -se na categoria difícil ( entre 15 \% e 40 \%, não incluso, de acertos ), sete são de nível médio ( entre 40 \% e 60 \%, não incluso, de acertos ), sete estão na categoria fácil ( entre 60 \% e 85 \%, não incluso, de acertos ) e uma teve a classificação muito fácil ( a partir de 85 \% de acertos ).
Nenhuma questão teve percentual de acertos menor do que 15 \% ( classificação muito difícil ).
Uma questão foi desconsiderada pela Comissão Assessora da Área de TECNOLOGIA EM DESIGN DE MODA.
A análise das questões objetivas do Componente Específico em TECNOLOGIA EM DESIGN DE MODA, em relação ao grau de facilidade da prova ( ver Tabela 3.8 ), mostrou que, das 27 questões, onze encontraram -se na categoria difícil ( entre 15 \% e 40 \%, não incluso, de acertos ), sete são de nível médio ( entre 40 \% e 60 \%, não incluso, de acertos ), sete estão na categoria fácil ( entre 60 \% e 85 \%, não incluso, de acertos ) e uma teve a classificação muito fácil ( a partir de 85 \% de acertos ).
Nenhuma questão teve percentual de acertos menor do que 15 \% ( classificação muito difícil ).
Uma questão foi desconsiderada pela Comissão Assessora da Área de TECNOLOGIA EM DESIGN DE MODA.
Das 27 questões objetivas relativas ao Componente Específico ( ver Tabela 3.9 ), três obtiveram o índice de discriminação muito bom, onze obtiveram o índice de discriminação bom e seis obtiveram nível de discriminação médio.
O índice fraco que eliminaria a questão foi alcançado por seis questões.
E, ainda, teve uma questão cancelada pela comissão.
Na sequência, no Quadro 3.2 estão especificadas, por questão, as habilidades aferidas pelas questões do Componente Específico da prova, bem como os resultados da análise relativa aos índices de facilidade e de discriminação de cada questão.
O Gráfico 3.8 mostra, a título de ilustração, o comportamento da questão 37 do Componente Específico para os estudantes concluintes da área de TECNOLOGIA DE MODA como função da nota geral na prova5.
Trata -se de uma questão considerada fácil para os concluintes de TECNOLOGIA EM DESIGN DE MODA.
O poder de discriminação foi bom para os estudantes da área.
Neste Gráfico, cada uma das cinco curvas, representa o percentual de respostas em determinada alternativa da questão, em função da nota total dos estudantes concluintes de TECNOLOGIA EM DESIGN DE MODA.
A curva roxa correspondeu à alternativa E, a correta para esta questão.
Assim, observou -se que entre os estudantes com escores mais baixos – desempenho geral mais fraco na prova como um todo – ocorreu a escolha de uma alternativa incorreta, a B – curva verde claro ou a alternativa D – curva verde escuro.
A proporção de estudantes que selecionaram a resposta correta E aumentou com o desempenho deste grupo de estudantes chegando ao acerto de 100 \% dos estudantes com notas pelo menos 5 Aqui também, as notas da prova foram grupadas em 8 intervalos de amplitude 12,5, fechados à direita e abertos à esquerda e um ponto com a nota máxima. 87,5.
Esta análise permitiu verificar como a questão discriminou entre grupos de desempenho de desempenho.
Os Gráficos relativos às demais questões do Componente Específico constam do Anexo I. 3.3 Análise das Questões Discursivas 3.3.1 Formação Geral A Tabela 3.10 apresenta as estatísticas básicas da prova por grupo de estudantes.
A análise dos resultados de desempenho dos estudantes nas questões discursivas que são apresentados na tabela 3.10 e no gráfico 3.9 mostrou que as notas foram mais baixas no conjunto dessas questões que no das questões objetivas.
O total de estudantes inscritos para a prova foi de 4.820.
Destes, 78,2 \% compareceram, sendo que a presença foi maior entre concluintes ( 87,1 \% ) do que entre ingressantes ( 74,3 \% ).
A média geral da prova foi de 27,80, sendo que os ingressantes obtiveram média mais baixa, de 26,73, que os concluintes, cuja média foi 29,85.
O desvio-padrão geral foi de 34,9, sendo o do grupo de ingressantes menor ( 34,4 ), em relação ao grupo dos concluintes ( 35,8 ), aqueles possuíram uma distribuição mais homogênea nas notas.
A nota máxima foi 100,0, atingida por concluintes e ingressantes.
Os coeficientes de assimetria para os dois grupos considerados, o de concluintes e o de ingressantes, bem como para a população como um todo foram todos positivos, indicando uma maior concentração à esquerda da distribuição.
Para todos os grupos um pouco mais da metade dos estudantes presentes obtiveram nota zero nas questões discursivas da Formação Geral.
Das 2.111 provas com nota zero, 30 ( 1,4 \% ) continham um protesto, 701 ( 33,2 \% ) estavam em branco e 1.380 ( 65,4 \% ) foram respondidas, porém de forma inadequada.
Entre os concluintes, a maioria ( 69,8 \% ) dentre aqueles com nota zero estavam com as provas respondidas de forma inadequada.
Entre os ingressantes, 63,2 \% das notas nulas decorreram de correção.
Uma comparação entre o desempenho de ingressantes e concluintes pode ser realizada a partir dos dados apresentados no Gráfico 3.9.
As notas dos concluintes, de um modo geral, foram ligeiramente mais elevadas do que as dos ingressantes, cuja distribuição de notas está à esquerda dos primeiros.
A ideia de melhor desempenho foi reforçada pela proporção de notas nulas : 53,4 \% entre os concluintes e 57,5 \% entre os ingressantes.
Na sequência, discorrer -se-á, especificamente, sobre uma das questões discursivas de Formação Geral, a número 10 já que a número 9 foi desconsiderada.
Em especial, serão analisados os conteúdos englobados na questão e os resultados obtidos pelos estudantes.
Ao final, são apresentados os comentários da banca corretora a respeito do observado nas respostas dos estudantes, suas impressões e conclusões.
Cumpre esclarecer que, tendo em vista que as questões discursivas de Formação Geral foram padronizadas, ou seja, constam de todas as provas, os comentários da banca foram os mesmos para todas as áreas acadêmicas, eis que direcionados a todos os estudantes que participaram do ENADE/2009.
Na sequência, discorrer -se-á, especificamente, sobre uma das questões discursivas de Formação Geral, a número 10 já que a número 9 foi desconsiderada.
Em especial, serão analisados os conteúdos englobados na questão e os resultados obtidos pelos estudantes.
Ao final, são apresentados os comentários da banca corretora a respeito do observado nas respostas dos estudantes, suas impressões e conclusões.
Cumpre esclarecer que, tendo em vista que as questões discursivas de Formação Geral foram padronizadas, ou seja, constam de todas as provas, os comentários da banca foram os mesmos para todas as áreas acadêmicas, eis que direcionados a todos os estudantes que participaram do ENADE/2009. 3.3.1.1 Análise da Questão Discursiva 9 de Formação Geral DESCONSIDERADA 3.3.1.2 Análise da Questão Discursiva 10 de Formação Geral A Questão Discursiva nº. 10 buscou avaliar as habilidades : • Ler e interpretar textos ; • Analisar e criticar informações ; • Extrair conclusões por indução e/ou dedução ; • Estabelecer relações, comparações e contrastes em diferentes situações ; • Fazer escolhas valorativas avaliando consequências ; • Questionar a realidade ; • Argumentar coerentemente.
O Quadro 3.3 apresenta as habilidades/competências, temas predominantes e estatísticas básicas da questão 10 de Formação Geral por grupo de estudantes.
Cumpre notar que como a questão 9 foi desconsiderada, a nota referente às questões discursivas da Formação Geral é idêntica à nota da questão 10.
A média geral desta questão foi 27,8, sendo que os ingressantes obtiveram média mais baixa, 26,7, que os concluintes, cuja média foi 29,8.
O desvio-padrão geral foi de 34,9, em relação ao grupo de ingressantes menor ( 34,4 ), indicando que, em contraposição ao grupo dos concluintes ( 35,8 ), aqueles possuíram uma distribuição mais homogênea nas notas.
A nota máxima foi 100,0, atingida por concluintes e ingressantes.
Os coeficientes de assimetria para os dois grupos considerados, o de concluintes e o de ingressantes, bem como para a população como um todo foram positivos, indicando uma maior concentração à esquerda da distribuição.
Para todos os grupos, cerca de metade dos estudantes presentes obtiveram nota zero nesta questão discursivas de Formação Geral.
Das 2.111 provas com nota zero, 30 ( 1,4 \% ) continham um protesto, 701 ( 33,2 \% ) estavam em branco e 1.380 ( 65,3 \% ) foram respondidas, porém de forma inadequada.
Entre os concluintes, a maioria ( 69,8 \% ) dentre aqueles com nota zero estava com as provas respondidas de forma inadequada.
Entre os ingressantes, 63,2 \% das notas nulas ( zero ) decorreram do desempenho do próprio estudante verificadas na correção.
Uma comparação entre o desempenho de ingressantes e concluintes para esta questão foi exatamente igual à já realizada para a nota discursiva de Formação Geral a partir dos dados apresentados no Gráfico 3.9. 3.3.1.3 Comentários sobre as respostas à Questão Discursiva 10 Apresentação A questão 10 de Formação Geral apresentou um questionamento capaz de analisar a habilidade do estudante em reconhecer o que é um direito social, optar por um direito social vigente que possa diminuir as desigualdades de renda familiar no Brasil e discutir os argumentos em favor do direito social escolhido como suporte à resposta.
Para análise das respostas apresentadas pelos estudantes, houve uma leitura inicial bastante criteriosa.
Foram considerados para efeito de pontuação, os seguintes argumentos : • Acesso à educação pública, gratuita e de qualidade, garantindo ao cidadão melhores condições de vida ; • Assistência à saúde, em seu contexto mais amplo, garantindo os cuidados essenciais de saúde, baseados em técnicas apropriadas ao alcance de todas as pessoas, contando com a participação da população ; • Permanência do estudante na escola, em todos os níveis escolares ( da educação infantil à educação superior, possibilitando a qualificação profissional ), o que garante melhores condições de trabalho e remuneração.
Considerou -se a qualificação e o desempenho docente, desenvolvendo nos estudantes melhores expectativas de aprendizagem e interesse pelo conteúdo da escola.
Os professores desenvolvendo metodologias adequadas e conhecimento pleno quanto a questões de aprendizagem ; • Condições dignas de trabalho, com remuneração que garanta qualidade de vida.
Foram consideradas as condições de trabalho, a preparação profissional para o trabalho e os cursos de aperfeiçoamento, oportunizando momentos de aprendizagem e desenvolvimento, a partir de situações desencadeadoras ; • Ser proprietário do imóvel em que reside, reduzindo os gastos com aluguel, garantindo melhores condições de vida, equilíbrio financeiro e construção do bem-estar familiar.
Desempenho dos estudantes na questão É importante salientar que havia respostas que desenvolveram adequadamente o questionamento apresentado, com fundamentação baseada em argumentos coerentes.
Contudo, em um número considerável de respostas, havia apenas um argumento e não dois, conforme fora solicitado.
Muitos estudantes informaram sobre problemas da educação, de um modo geral.
Algumas \textbf{observações} mencionadas pelos corretores durante o processo de correção foram : • há preocupação com uma educação de qualidade, com sugestões de cursos profissionalizantes e métodos preparatórios de aulas. • a educação é revelada como de baixa qualidade, exaltando -se treinamentos e cursos de atualização para a classe docente. • a saúde como essência para a vida.
O SUS é colocado como parte integrante de direito social, garantindo às pessoas fácil acesso. • oportunidade de emprego como forma de ascensão social, direito, cidadania. • os políticos brasileiros como causadores do despreparo educacional.
As respostas analisadas apresentaram norma culta, coesão e coerência razoáveis, que garantiram a compreensão da mensagem desenvolvida pelos estudantes.
Os erros ortográficos foram poucos e não comprometeram a discussão sobre o assunto e sobre os pontos de vista abordados.
Quanto à coerência, houve um desvio quanto ao objetivo da questão, já que alguns estudantes, através dos seus argumentos retrataram reivindicações, reclamações e situações de problemas enfrentados pelos brasileiros, citando golpes políticos.
Houve erros quanto à concordância verbal e nominal.
Houve também, dificuldade em utilizar os conectores corretamente.
Alguns estudantes, por falta de ideias, apresentaram em suas respostas redundância e prolixidade.
O desempenho foi considerado satisfatório, pois a maioria dos estudantes soube responder ao questionamento.
Houve domínio do conteúdo, pois as ideias, os argumentos e a originalidade corresponderam ao esperado. 3.3.2 Componente Específico Nas questões discursivas do Componente Específico – questões de números 38, 39 e 40 –, a diferença de desempenho entre os grupos não foi muito acentuada, inclusive menor do que para as questões discursivas de Formação Geral.
A Tabela 3.11 contém as estatísticas da média destas questões.
Observa -se que no grupo dos concluintes, a nota média alcançada foi 16,38, enquanto que no grupo dos ingressantes foi de 9,25, ambas menores do que as obtidas nas questões discursivas de Formação Geral, que foram de 26,7 para os ingressantes e 29,8 para os concluintes ( Tabela 3.10 ).
As médias foram também menores dos que as obtidas para as questões objetivas do Componente Específico ( Tabela 3.7 ) : respectivamente 51,5 e 59,5 para ingressantes e concluintes.
A análise de cada uma destas questões será feita nas subseções seguintes.
Os valores positivos para a medida de assimetria indicaram uma distribuição com cauda para a direita.
Dos 4.820 presentes, 36,3 \% obtiveram nota zero.
A tabela 3.12 apresenta um detalhamento da situação destes estudantes.
A proporção dos que responderam a todas as questões correspondeu aproximadamente a 12 \% das provas.
Não houve prova toda em branco, nem toda com protesto, nem combinação das duas modalidades.
O Gráfico 3.10 apresenta a distribuição das notas nas questões discursivas do Componente Específico por grupo de estudantes.
A moda se situou no primeiro intervalo, aquele com valores entre 0 e 10 incluindo os extremos do intervalo.
A impressão de assimetria informada pelos coeficientes de assimetria positivos ( Tabela 3.11 ) foi reforçada no gráfico pela concentração no primeiro intervalo.
O Gráfico 3.10 apresenta a distribuição das notas nas questões discursivas do Componente Específico por grupo de estudantes.
A moda se situou no primeiro intervalo, aquele com valores entre 0 e 10 incluindo os extremos do intervalo.
A impressão de assimetria informada pelos coeficientes de assimetria positivos ( Tabela 3.11 ) foi reforçada no gráfico pela concentração no primeiro intervalo. 3.3.2.1 Análise da Questão Discursiva 38 do Componente Específico A questão 38, cujas habilidades aferidas encontram -se descritas no Quadro 3.4, que ainda contém outras informações, seguido da Tabela 3.13 e do Gráfico 3.11, foi aquela em que os grupos de estudantes observados obtiveram um desempenho intermediário entre as questões discursivas do Componente Específico.
A nota média do grupo dos ingressantes foi de 12,68, e a dos concluintes 23,30.
O percentual de notas zero obtidas nos dois grupos foi de 55,8 \%, maior entre os ingressantes ( 62,7 \% ) do que entre os concluintes ( 42,4 \% ).
Mais de 80 \% dos concluintes e dos ingressantes que tiraram zero deixaram a questão em branco ( ver Tabela 3.13 ).
Uma pequena minoria ( 6,8 \% ) utilizou a questão para protestar.
Apesar do baixo desempenho, em ambos os grupos a nota máxima de 100,0 pontos foi alcançada.
Os coeficientes de assimetria positivos indicam uma inclinação para a esquerda e cauda para a direita.
O Gráfico 3.11 apresenta a distribuição das notas na questão 38 discursiva do Componente Específico por grupo de estudantes.
A moda se situou no primeiro intervalo, aquele com valores entre 0 e 10 incluindo os extremos do intervalo.
A impressão de assimetria informada pelos coeficientes de assimetria positivos ( Quadro 3.4 ) foi reforçada no gráfico pela concentração no primeiro intervalo. 3.3.2.2 Comentários sobre as respostas à Questão Discursiva 38 Essa questão foi elaborada tendo como conteúdo predominante a história da indumentária da moda e do design, focada no século XX, a partir de criadores como Kenzo, seguido por Issey Miyake, Rei Kawakubo e Yohji Yamamoto.
A especificação da questão para descrição foi dividida em A, B, C para justificar os diferentes estilos no contexto histórico.
O padrão de respostas ofereceu opções claras para a correção, o que facilitou a ação dos corretores que concluíram que o item C apresentou maior dificuldade para os estudantes.
Considerando que o assunto da questão ofereceu uma abordagem bem contemporânea sobre a moda, não se justifica que grande número de estudantes deixasse essa questão em branco.
Os estudantes que responderam tal questão fixaram seu empenho nos itens A, B e muitos tiveram conclusões de acordo com o que foi solicitado, mesmo porque apresentar argumentos justificando a entrada dos japoneses na moda parisiense é fato evidente, conhecido e divulgado em quase todas as bibliografias vigentes.
As redações das respostas apresentavam alguma deficiência na ortografia, mas considerando o texto como um todo, a coerência e a coesão foram evidenciadas. 3.3.2.3 Comentários sobre as respostas à Questão Discursiva 38 Essa questão foi elaborada tendo como conteúdo predominante a história da indumentária da moda e do design, focada no século XX, a partir de criadores como Kenzo, seguido por Issey Miyake, Rei Kawakubo e Yohji Yamamoto.
A especificação da questão para descrição foi dividida em A, B, C para justificar os diferentes estilos no contexto histórico. 3.3.2.3 Análise da Questão Discursiva 39 do Componente Específico DESCONSIDERADA 3.3.2.4 Análise da Questão Discursiva 40 do Componente Específico A questão 40, cujas habilidades a serem aferidas encontram -se descritas no Quadro 3.5, que ainda contém outras informações, seguido da Tabela 3.14 e do Gráfico 3.12.
A questão foi aquela em que os grupos de estudantes estudados obtiveram o melhor desempenho entre as questões discursivas do Componente Específico.
A nota média do grupo dos ingressantes foi de 15,13, e a dos concluintes 25,90.
Nesta questão ocorre também o menor percentual de notas nulas : 53,6 \% entre os ingressantes e 30,1 \% entre os concluintes.
Aqui também a maioria das provas com nota zero decorre de provas deixadas em branco : 69,6 \% entre os ingressantes e 66,8 \% entre os concluintes ( ver Tabela 3.14 ).
Uma pequena minoria ( 3,6 \% ), dentre aqueles com nota nula, utilizou a questão para protestar.
Em ambos os grupos a nota máxima de 100,0 pontos foi alcançada.
Os coeficientes de assimetria positivos indicam uma inclinação para a esquerda com uma cauda para a direita.
O Gráfico 3.12 apresenta a distribuição das notas na questão 40 discursiva do Componente Específico por grupo de estudantes.
A moda se situou no primeiro intervalo, aquele com valores entre 0 e 10 incluindo os extremos do intervalo.
A impressão de assimetria informada pelos coeficientes de assimetria positivos ( Quadro 3.5 ) é reforçada no gráfico pela concentração neste primeiro intervalo.
Existem ainda outros máximos locais.
O Gráfico 3.12 apresenta a distribuição das notas na questão 40 discursiva do Componente Específico por grupo de estudantes.
A moda se situou no primeiro intervalo, aquele com valores entre 0 e 10 incluindo os extremos do intervalo.
A impressão de assimetria informada pelos coeficientes de assimetria positivos ( Quadro 3.5 ) é reforçada no gráfico pela concentração neste primeiro intervalo.
Existem ainda outros máximos locais. 3.3.2.5 Comentários sobre as respostas à Questão Discursiva 40 A questão apresenta uma situação problema bem atual, focada na Tecnologia da Confecção, buscando a habilidade e a competência pela elaboração da ficha técnica para orientação da produção seriada.
O enunciado da questão solicitou uma descrição de etapas de uma confecção de aventais ( A ) com indicação de pelo menos dois fatores limitadores neste empreendimento ( B ).
Portanto, as duas respostas eram dependentes e seqüenciadas.
A grande maioria dos estudantes deixou toda a questão em branco e os que a responderam, muitas vezes apresentaram as etapas sem descrevê -las.
Alguns poucos estudantes, vagamente indicaram os fatores limitantes da utilização dos aventais, não respondendo o item B.
A partir desta explicitação, concluiu -se que o conteúdo focado na questão, embora muito atual, não deve ter sido visto na grande maioria dos cursos, daí a razão das questões em branco ou incompletas.

% matched lemmas: observação
\end{textsample}
